% !TeX root = ../navier-stokes-manuscript.tex
\subsection{The pressure-complete mixed jet}\label{sub:ThePressureCompleteMixedDerivativeTower}

There is no separate spatial history to add to the temporal Tower. Differentiate the same rows in space. For \(n\in\Nzero\) and \(\alpha\in\Nzero^3\), define
\begin{equation}\label{eq:BodyMixedJetDefinition}
  J_{n,\alpha}
  :=\partial_t^n\partial_x^\alpha u,
  \qquad
  \Pi_{n,\alpha}
  :=\partial_t^n\partial_x^\alpha p.
\end{equation}
The two indices remain independent. Increasing \(n\) follows the fixed-coordinate response in time. Increasing \(|\alpha|\) resolves finer spatial structure inside that chosen response.

Apply \(\partial_x^\alpha\) to the temporal recurrence \eqref{eq:BodyTemporalVelocityRecurrence}.
For \(\beta\leq\alpha\), write
\[
  \binom{\alpha}{\beta}
  :=\prod_{j=1}^3\binom{\alpha_j}{\beta_j}.
\]
The temporal and spatial product rules give
\begin{align}
  J_{n+1,\alpha}
  &+
  \sum_{a=0}^{n}
  \sum_{\beta\leq\alpha}
  \binom{n}{a}\binom{\alpha}{\beta}
  (J_{a,\beta}\cdot\nabla)
  J_{n-a,\alpha-\beta}
  +\nabla\Pi_{n,\alpha}
  \notag\\
  &=\nu\Delta J_{n,\alpha},
  \qquad
  \nabla\cdot J_{n,\alpha}=0.
  \label{eq:BodyMixedJetRecurrence}
\end{align}
The pressure coordinate is fixed by the same velocity rows:
\begin{equation}\label{eq:BodyMixedPressureRecurrence}
  -\Delta\Pi_{n,\alpha}
  =\partial_x^\alpha\partial_i\partial_j
  \sum_{a=0}^{n}\binom{n}{a}(V_a)_i(V_{n-a})_j.
\end{equation}
The Riesz normalization in \eqref{eq:PressureNormalization} chooses the pressure row in the decaying whole-space class. Every term still belongs to the original history. The binomial coefficients record how time and space derivatives divide between the transporting and transported factors; viscosity reaches two spatial derivatives above the present coordinate; pressure returns the nonlocal response imposed by incompressibility.

The complete mixed jet is the family
\[
  \mathcal J[u_0]
  :=\left\{
    J_{n,\alpha},\Pi_{n,\alpha}:
    n\in\Nzero,\ \alpha\in\Nzero^3
  \right\}
\]
together with the recurrences and divergence relations above. These are derivatives of one \((u,p)\), not an infinite norm and not a collection of independently chosen fields. At \(t=0\), the same recurrences determine every finite row from \(u_0\); Appendix~\ref{app:DerivativeAndFiniteRecordCalculus} writes that initialization in full.

The indices now tell us where every finite derivative sits. They do not yet tell us which of those coordinates the equation can control as \(t\uparrow T_*\). For that, return to critical height and differentiate its viscous balance.
